\section{Marco teorico}

\subsection{Redes P2P}

En una red peer-to-peer cada nodo puede actuar como cliente y servidor. Esta
propiedad permite que un cliente solicite datos a otros clientes sin depender
de un unico repositorio de archivos. El proyecto usa un servidor central solo
para registrar metadatos y sembrar vecinos iniciales; las transferencias de
datos y la busqueda distribuida ocurren directamente entre clientes.

\subsection{Calidad de servicio}

La calidad de servicio se refleja en tiempos de busqueda, disponibilidad de
pares y continuidad de las transferencias. La implementacion reduce bloqueos
separando responsabilidades en hilos: la consola queda en el hilo principal,
el puerto de datos tiene un listener dedicado y cada transferencia usa hilos de
segmento. Tambien se ignora \texttt{SIGPIPE} y se revisan errores de I/O para
evitar que una conexion cerrada termine el proceso.

\subsection{Hashing de archivos}

El sistema identifica un archivo por el par \texttt{(hash, size)}. Se usa una
funcion propia basada en FNV-1a de 64 bits sobre el contenido del archivo, lo
que permite reconocer el mismo archivo aunque cambie su nombre. El tamano se
incluye como segunda llave para disminuir falsos positivos y para dividir los
rangos de descarga.

\subsection{Sockets TCP y framing}

Todos los mensajes viajan sobre TCP. Antes de cada mensaje se envia una
longitud externa con \texttt{net\_send\_msg}; dentro de ese frame aparece un
\texttt{p2p\_msg\_header\_t} con opcode, version y tamano del payload. Los
enteros de multiples bytes se codifican en orden de red con \texttt{htons},
\texttt{htonl} y conversiones de 64 bits equivalentes.

\subsection{Busqueda distribuida}

La busqueda distribuida usa flooding con TTL. Cada consulta incluye un ID de
16 caracteres hexadecimales, termino de busqueda, IP y puerto del originador,
y un contador TTL. Cada cliente mantiene una cache de IDs vistos durante 60
segundos para evitar ciclos. Cuando recibe una consulta nueva, revisa su
carpeta compartida, responde directamente al originador si encuentra
coincidencias y reenvia a sus vecinos con TTL decrementado.

El TTL por defecto es 4 y el maximo aceptado es 6. Con grado promedio 3, un TTL
4 tiene un limite teorico de aproximadamente $3^4 = 81$ mensajes, suficiente
para una prueba local de tres clientes sin saturar la LAN. La ventana de
agregacion por defecto es 2 segundos y puede configurarse al iniciar el
cliente.
